\section{Mathematical taste from equality cases}\label{ch09:sec:taste}
Many arguments begin with an inequality, and the inductive proof of Hall's theorem uses one inequality in two different ways. When $|N(S)|\ge|S|$ holds with equality, the group $S$ is rigid: any full assignment must use up all of $N(S)$ on the students of $S$. That rigidity told us how to split the problem into two independent pieces. When the inequality is strict, whole numbers give a margin of at least one, and that margin was enough to absorb the loss of one task. The proof was found by asking separately what each situation allows.

The rigidity in the equality case is a familiar fact in disguise. Let $f:A\to B$ be an injective function between finite sets of the same size. Its $|A|$ outputs are different members of $B$, and $B$ has only $|A|$ members, so every member of $B$ is an output: $f$ is also surjective. The tight-block argument is exactly this, with the students of $S$ as $A$, the tasks of $N(S)$ as $B$, and the matching as $f$.

The same two questions are worth asking whenever you meet a bound. Besides asking whether the bound can be improved, ask what happens when it is attained. If equality forces a recognizable structure, that structure may let you break the problem into smaller pieces. If the inequality is strict, ask how large the margin is and what it can pay for. You will meet a margin of this second kind in Chapter~\ref{ch:contraction}. The contraction theorem (Theorem~\ref{thm:banach}) requires a number $q$ with $0\le q<1$, and the positive gap $1-q$ is what keeps its error bounds finite.
